\parttitle{Mathematical destination}
\labelled{Three possible outcomes.}{With exactly one fixed heavy label among nine and equal numbers placed on each pan, two simulated weighings suffice and one cannot suffice. A test returns LEFT, RIGHT or BALANCED. Split the candidates into three equal groups; its answer selects one group. The model assumes truthful feedback and a fixed secret. It says nothing about real objects of unknown or unequal masses.}
\labelled{Different objectives.}{For letters A, B, C, D occurring 4, 2, 1, 1 times, the least total number of truthful yes/no questions over the eight messages is 14. One optimal plan has depths 1, 2, 3, 3 and worst case 3; a balanced two-question plan costs 16. Identification ends when one candidate remains. Questions can ask about any subset. The same decision plan is used for each message.}
\labelled{Words need boundaries.}{Distinct variable-length words need not give distinct concatenated messages. Book 1 has a collision: B and AC both give RY. Book 2 is prefix-free and uniquely decodable. Book 3 is also uniquely decodable although A's word is a prefix of B's word; prefix freedom is sufficient, not necessary. Only rows made by the book need be decodable. An arbitrary R/Y row may be invalid.}
\labelled{What children discover.}{A new kind of test changes the information available; a frequent message deserves a short route; removing spaces can lose information. Trials establish examples. A guarantee needs a plan for every outcome, a cost comparison needs all possible tree shapes, and a uniqueness explanation needs a decoding argument.}
\labelled{Band map.}{Student pages 1 and 3 invite grades 2--5; page 2 invites grades 4--5. K--1 can enact an adult-led three-object weighing game, without being expected to optimize nine-object searches or read variable words independently. Counting up to nine and sorting objects suffice for weighing. Message words require following left-to-right order; page 2 adds totaling eight small counts and comparing adaptive plans. No probability, logarithms or formal trees are entry prerequisites.}
\labelled{Status.}{Prepared, unpiloted return-visit material. These three investigations share Week 6's information theme but are separate from its fixed-length counter scores and worst-case yes/no questions. They are not a single-hour completion quota.}
\newpage
\parttitle{Materials, first route and launch}
\labelled{Per pair.}{Nine numbered slips 1--9, two paper pan areas, one folder, a secret-label slip, pencils and scraps. For page 2, four distinguishable objects labelled A--D and eight slips A,A,A,A,B,B,C,D. For page 3, about 20 two-colour counters or R/Y tiles and letter slips. Real weights or a balance are unnecessary: ``heavy'' means the keeper's secret label. The pan boxes are approximately 3.05 by 1.35 inches; slips should be less than 0.7 inches across. Print US Letter at 100\%. Physical fit and this procedure have not been rehearsed.}
\labelled{Prepare.}{Make one nine-label set per pair. The keeper writes the fixed secret behind the folder and returns LEFT or RIGHT if that pan contains its card, otherwise BALANCED. Check equal pan counts before giving feedback. After identification reveal the slip and replay the weighings to check them. For the upper trio rotate breaker, keeper and referee.}
\labelled{Launch together, 3--5 minutes.}{Let children sort and handle the slips. With practice labels 10,11,12, the keeper secretly picks 12, puts 10 on the left and 12 on the right, and points RIGHT. Show another test with 12 off both pans and say BALANCED. Ask the group to make a legal test; then distribute different starting pages. The printed non-task visual preserves the weighing convention.}
\labelled{A manageable visit.}{Spend 15--25 minutes on page 1, with partners exchanging roles. Save a plan and a hard secret. A later 20--30-minute visit may start with the physical batch on page 2 or word tiles on page 3. A younger table can play the three-object search and short joined-message collisions while an adult reads or scribes. Use the current three anchored adults at KK11 / 3333 / 445 tables; no child needs to finish every page.}
\labelled{Facilitation.}{Let children choose groups, question order and message splits. Keep tested plans visible. An adult may read labels, operate feedback or scribe a tree; avoid choosing all the tests for a child. Offer the hints on later pages only after genuine attempts.}
\labelled{Lineage and evidence.}{Current Week 6 F06-v4 covers fixed-row match scores and worst-case questions; current Week 18 already covers damaged fixed-length transmissions, deliberately excluded here. These task constructions and proofs are original derivations. \emph{Math Circle by the Bay}, preface printed pp. viii--x (PDF pp. 9--11), actually describes manipulatives, reserve challenges, explanations and differing depth. Our staffing, timings and companion format are adaptations, not arrangements reported there.}
\newpage
\parttitle{Problem 1 / A different kind of test}
\labelled{Checked answer.}{Two weighings are necessary and sufficient. First put 1,2,3 on LEFT and 4,5,6 on RIGHT. LEFT leaves candidates 1,2,3; RIGHT leaves 4,5,6; BALANCED leaves 7,8,9. Within the surviving group put one candidate on each pan and leave the third off. LEFT, RIGHT and BALANCED identify the respective candidate. The second comparison uses equal pan sizes.}
\labelled{Why fewer cannot work.}{One weighing has at most three responses. Each response must account for all the secrets producing it. Nine different secrets cannot all produce different responses. No clever choice of pan groups can create more than three answers. The information argument applies to any one weighing, including one that leaves many objects off.}
\labelled{How to check a child's plan.}{For every possible secret, run the keeper's responses without changing the secret. Check whether the stated outcome branch ends at one label. A plan that works for one friendly keeper is an experiment; a plan with a tested branch for all nine secrets is a guarantee. The breaker need not weigh a final candidate merely to confirm it.}
\labelled{Entry and questions.}{Begin with three labelled objects and one comparison. For nine, ask which labels are still possible after a response. If groups are unequal, let children see whether the largest surviving group can be finished with one more comparison. Ask how many labels could be handled with three comparisons.}
\labelled{Hint if needed.}{A BALANCED outcome should tell you something too. Lay remaining candidates in three piles before deciding which two piles to put on the pans.}
\labelled{Continuation.}{With 27 labels, three successive splits into thirds work; fewer than three tests cannot guarantee it because two tests have at most nine outcome sequences. Arbitrary group sizes need care: an equal-pan constraint may forbid a desired split, though uniform powers of three avoid that issue. Do not generalize this result to a single odd object that could be either heavier or lighter.}
\labelled{Record for a return visit.}{Keep the actual pan groups and all outcome branches, noting which secrets were physically played and which branches were reasoned through. No observations about children's understanding are available yet.}
\newpage
\parttitle{Problem 2 / Cheapest whole batch}
\labelled{Checked answer.}{Ask ``Is it A?'' If yes, stop. If no, ask ``Is it B?'' If yes, stop. If no, ask ``Is it C?'' A yes identifies C; a no identifies D. Costs are A:1, B:2, C:3, D:3. Across the printed batch the cost is $4\cdot1+2\cdot2+1\cdot3+1\cdot3=14$. The most any one message needs is three questions. Swapping C and D also works.}
\labelled{A contrasting plan.}{Ask whether the letter belongs to \{A,B\}, then distinguish the remaining pair. Every letter needs two questions; eight messages cost 16. It wins the worst-case contest but loses the batch-cost contest. The printed task permits adaptivity within a message, but requires the same plan for all messages; it does not allow exploiting which batch cards have already appeared.}
\labelled{Why 14 is minimal.}{Remove questions with only one possible response. A useful four-leaf binary decision tree either divides its letters 2-and-2 at the root, giving depths 2,2,2,2, or divides 1-and-3, giving depths 1,2,3,3. In the latter shape it is cheapest to assign shorter routes to more frequent letters: exchanging frequencies $f>g$ at depths $d>e$ saves $(f-g)(d-e)>0$. Thus A gets depth 1, B gets depth 2 and C,D get depth 3. The two possibilities cost at least 16 and 14 respectively. No probability notation is needed.}
\labelled{Entry and questions.}{Act out a plan eight times and put a counter in a common cost bowl after each question. Compare bowls for two plans. Ask whether two plans with the same worst case can have different total cost. If drawing trees distracts children, arrange physical question slips instead.}
\labelled{Hint if needed.}{Which message occurs most often? Try making its route short, then see what happens to the others. Avoid telling children the complete order.}
\labelled{Continuation.}{Invent a different frequency batch. All equal frequencies favour a balanced tree; sufficiently skewed frequencies favour an unequal one. Children can trade physical plan cards and test total costs. A broad optimal-coding algorithm is a later direction, not a prerequisite.}
\labelled{Record.}{Save the plan, its four route lengths and total. Note whether a child independently distinguished worst-case and total cost, or an adult supplied that distinction.}
\newpage
\parttitle{Problem 3 / Where did the spaces go?}
\labelled{Book 1.}{A=R, B=RY, C=Y. B and AC are different messages but both produce RY. Longer collisions follow by adding identical prefixes or suffixes to these messages. One collision proves that this book is ambiguous; finding many is optional.}
\labelled{Book 2.}{A=R, B=YR, C=YY. Every nonempty valid row can be decoded from the left: an initial R must be A; an initial Y must belong to a two-counter word, whose second symbol decides B or C. Remove that word and repeat. This explains uniqueness for messages of every length, not just the examples tried. A single final Y is invalid rather than ambiguous.}
\labelled{Book 3.}{A=R, B=RY. Every valid row is uniquely decodable. Read from the right: a last R must be A; a last Y must be the end of B, whose preceding R belongs to that same B. Remove the last word and repeat. Equivalently, each Y must attach to the R just before it. This book demonstrates why ``one word starts another'' does not by itself prove ambiguity. A row containing YY cannot be a valid message.}
\labelled{Printed example.}{Practice words M=YR and N=R turn MN into YR\,R, then YRR. The divider in the middle picture marks the word boundary; the last picture removes it. These are different practice letters, not a solution to one of the three books.}
\labelled{Entry and questions.}{Make each word with counters on a slip. Place two or three letter slips in order, concatenate their counter rows, and hide the slips. The receiver may try moving paper boundary markers. Ask whether ``I tried many and found no collision'' is enough to establish uniqueness. An adult can write the decoded letters while the child chooses the boundaries.}
\labelled{Hint if needed.}{For Book 1 compare one long word with two short words. For Books 2 and 3 ask which word must be at one end of a valid row. The useful end need not always be the left end.}
\labelled{Continuation and limits.}{Children can invent a fourth book and ask a partner to find a collision or a decoding argument. Prefix-free books have immediate left-to-right decoding, but unique decoding can depend on a later symbol or reading from the right. The task assumes the receiver knows the book and the row's start and end; unknown book choice is a different problem.}
\labelled{Record.}{Save a collision pair or the child's end-reading argument, and mark which book was actually explored. This is a segmentation investigation, not a damaged-counter correction exercise.}
